\documentclass{article}

\usepackage[margin=1in]{geometry}
\usepackage{booktabs}
\usepackage{tabularx}
\usepackage{longtable}
\usepackage{array}
\usepackage{pdflscape}
\usepackage{enumitem}
\usepackage{hyperref}

\hypersetup{
    colorlinks=true,       % false: boxed links; true: colored links
    linkcolor=red,          % color of internal links (change box color with linkbordercolor)
    citecolor=green,        % color of links to bibliography
    filecolor=magenta,      % color of file links
    urlcolor=cyan           % color of external links
}

\setlength{\emergencystretch}{3em}

\title{Hazard Analysis\\\progname}

\author{\authname}

\date{}

\input{../Comments.text}
\input{../Common.text}

\newcommand{\srs}{\href{../SRS-Volere/SRS.pdf}{SRS}}

\begin{document}

\maketitle
\thispagestyle{empty}

~\newpage

\pagenumbering{roman}

\begin{table}[hp]
\caption{Revision History} \label{TblRevisionHistory}
\begin{tabularx}{\textwidth}{llX}
\toprule
\textbf{Date} & \textbf{Developer(s)} & \textbf{Change}\\
\midrule
Oct 19, 2026 & William Pol & First draft of all sections; see the
disclosure appendix\\
\bottomrule
\end{tabularx}
\end{table}

~\newpage

\tableofcontents

~\newpage

\pagenumbering{arabic}

\section{Introduction}
\label{sec:intro}

This document analyzes the hazards of \progname{}, which predicts a target
model's key-value (KV) cache from a source model's KV cache and loads the
prediction into a large language model (LLM) serving engine. When a
semantic router sends the next turn of a conversation to a different model,
\progname{} lets the new model start from the predicted cache instead of
processing the whole conversation history again (prefill). The prediction
is approximate, so much of this analysis concerns how a wrong prediction is
caught before it reaches the user. The Software Requirements
Specification (\srs{}) describes the requirements this analysis refers to.

A \emph{hazard} is a condition of \progname{} that, together with a
condition in its environment, can lead to a loss. \progname{} does not
control physical equipment, so, as the course allows, we treat harm
broadly: incorrect answers, disclosure of private information, unavailable
service, wasted computing resources, and misleading performance claims.

\paragraph{Changes to the template.} We added a hazard list
(Section~\ref{sec:hazards}) before the Failure Mode and Effect Analysis
(FMEA) so that the table, the \srs{} and the roadmap can all refer to the
same hazard identifiers. The FMEA table adds a hazard column for that
traceability.

\section{Scope and Purpose of Hazard Analysis}
\label{sec:scope}

The purpose of this analysis is to find how \progname{} can fail, decide
how each failure is detected and mitigated, and turn the mitigations into
requirements. Following the as low as reasonably practicable (ALARP)
principle, the goal is not zero risk but risk reduced as far as is
practical within the capstone.

The losses considered are:
\begin{enumerate}[label=L\arabic*., ref=L\arabic*]
  \item \label{loss:wrong} \textbf{Incorrect response.} An application
  receives an answer that differs from what the selected model would have
  produced, and nothing marks it as unusual.
  \item \label{loss:privacy} \textbf{Privacy loss.} One user's
  conversation, instructions or answer is revealed to another user,
  directly or indirectly.
  \item \label{loss:service} \textbf{Lost service.} A valid request
  fails, times out, or slows other requests beyond acceptable limits.
  \item \label{loss:waste} \textbf{Wasted resources.} Graphics processing
  unit (GPU) memory or time is spent on work that is never used.
  \item \label{loss:claims} \textbf{Misleading results.} Evaluation reports
  claim a speedup or quality level that the system does not deliver,
  leading stakeholders to adopt it wrongly.
\end{enumerate}

The analysis covers \progname{} on a single server, as the \srs{} scopes
it, including the development and evaluation of the KV cache predictor. It
excludes the semantic router's choice of model, the application that
stores the conversation, training of the base models, and operation across
several servers, because these are outside the product boundary in the
\srs{}. Failures at those boundaries are still considered where they
affect \progname{}.

\section{System Boundaries and Components}
\label{sec:components}

The components below come from the context and work partitioning in the
\srs{}. They are used to brainstorm hazards, not as a design.

\paragraph{Inside \progname{}.}
\begin{enumerate}[label=C\arabic*., ref=C\arabic*]
  \item \label{comp:validate} \textbf{Request validation:} checks request
  fields, sizes and settings before work begins.
  \item \label{comp:predict} \textbf{KV cache predictor:} estimates the
  target model's keys and values from the source model's KV cache, using
  versioned parameters and recorded token, layer and head alignment
  (\srs{} FR-16). Its parameters may be fitted on paired source and target
  caches.
  \item \label{comp:check} \textbf{Compatibility checker:} decides whether a
  candidate cache may be used, using the checks in the \srs{} formal
  specification (model setup, exact prefix, permission, readiness and
  resources).
  \item \label{comp:load} \textbf{KV loader:} places an accepted cache into
  the serving engine's KV memory and prefills the remaining prompt.
  \item \label{comp:native} \textbf{Native prefill fallback:} processes the
  full prompt normally when reuse is off, unavailable or unsuitable.
  \item \label{comp:output} \textbf{Decoding and output streaming:} returns
  generated tokens in order under the correct request identifier (ID).
  \item \label{comp:life} \textbf{Request lifecycle:} handles completion,
  failure, cancellation and timeout, and releases memory.
  \item \label{comp:config} \textbf{Configuration and tested-setup
  registry:} stores the approved model/cache setups and the reuse on/off
  control.
  \item \label{comp:log} \textbf{Logging and metrics:} records per-request
  status, reuse decisions and measurements.
  \item \label{comp:bench} \textbf{Benchmark tools:} run matched workloads
  with and without \progname{} and produce comparison reports.
\end{enumerate}

\paragraph{Outside \progname{}, but sources of hazards.}
\begin{enumerate}[resume*]
  \item \label{comp:router} \textbf{Semantic router:} selects the target
  model for each turn.
  \item \label{comp:app} \textbf{Application:} supplies the conversation,
  the new message and the sharing group of each request.
  \item \label{comp:engine} \textbf{Serving engine and GPU:} the vLLM
  engine \cite{Kwon2023vLLM}, its memory manager, the source and target
  model files, the source model's KV cache and the hardware.
\end{enumerate}

\section{Critical Assumptions}
\label{sec:assumptions}

The \srs{} lists assumptions A1--A6. Only the following affect which
hazards are analyzed:

\begin{itemize}
  \item \textbf{A1 (trusted cache provider).} The \srs{} assumes the
  initial cache provider runs in the operator's trusted environment and
  that users cannot upload caches. Because \progname{}'s own predictor now
  produces the cache, deliberate tampering is less likely, but a cache can
  still be corrupted by a software fault, a storage error or a mismatched
  predictor version. We therefore do \emph{not} rely on A1: H9 remains in
  the analysis, and stronger mitigations are placed on the roadmap for any
  future release that accepts caches from outside.
  \item \textbf{A4 (a useful predictor exists).} The \srs{} states that a
  predictor which is both accurate enough and cheaper than prefill is a
  proof-of-concept question. We do not assume it succeeds: H8 and H15
  cover a predictor that is inaccurate or too slow, and in both cases the
  native prefill path must remain correct.
  \item \textbf{A2 (router and application inputs).} The router's model
  choice and the application's conversation are taken as given. \progname{}
  still validates them (H1, H6), because a wrong input must not cause a
  wrong cache to be used.
  \item \textbf{A3 (adequate server).} The server can run the selected
  models normally. Running out of memory is still analyzed (H6, H12, H14).
  \item \textbf{A5 (data for the predictor).} Public or synthetic
  conversations are available for fitting and testing the predictor. If
  the same conversations are used for both, results look better than they
  are; this is analyzed as part of H11.
\end{itemize}

No assumption states that a component never fails.

\section{Hazards}
\label{sec:hazards}

Table~\ref{tab:hazards} lists the hazards.

\begin{longtable}{@{}p{0.07\textwidth}p{0.62\textwidth}p{0.22\textwidth}@{}}
\caption{Hazards and the losses they can cause}\label{tab:hazards}\\
\toprule ID & Hazard & Losses\\\midrule\endfirsthead
\toprule ID & Hazard & Losses\\\midrule\endhead
H1 & An incompatible cache (wrong model, tokenizer, format, positions or
precision) is used and changes the response. & \ref{loss:wrong}\\
H2 & One request uses or reveals another request's private history,
instructions or answer. & \ref{loss:privacy}\\
H3 & An old or incomplete cache is used. & \ref{loss:wrong}\\
H4 & Cache data is changed or removed while a request still uses it. &
\ref{loss:wrong}, \ref{loss:service}\\
H5 & Logs or measurement files reveal private data. &
\ref{loss:privacy}\\
H6 & Memory or other resources run out. & \ref{loss:service}\\
H7 & Output from a failed reuse attempt is mixed with a restarted answer. &
\ref{loss:wrong}\\
H8 & A compatible, tested predicted cache still lowers answer quality,
because cross-model prediction is approximate. & \ref{loss:wrong},
\ref{loss:claims}\\
H9 & A cache is corrupted or deliberately altered before it is loaded. &
\ref{loss:wrong}, \ref{loss:privacy}\\
H10 & The time difference between reuse and normal processing reveals
information about other users' conversations. & \ref{loss:privacy}\\
H11 & Evaluation reports overstate the benefit or hide failures. &
\ref{loss:claims}\\
H12 & Predicted caches that are never used consume memory and GPU time. &
\ref{loss:waste}, \ref{loss:service}\\
H13 & The model, tokenizer, alignment rules or predictor version change
between predicting a cache and using it. & \ref{loss:wrong}\\
H14 & Many reuse attempts fail at once and the fallback overloads the
server. & \ref{loss:service}\\
H15 & Predicting, checking and loading the cache takes longer than native
prefill, so reuse slows the request instead of speeding it up. &
\ref{loss:service}, \ref{loss:waste}\\
\bottomrule
\end{longtable}

Three hazards need context.

\paragraph{H8, approximate prediction.} Unlike same-model prefix caching,
predicting one model's KV cache from another's does not reproduce the
target model's own cache exactly. Published results show that the quality
loss can be large and depends on the direction of transfer. For example,
one cross-family method reports a question-answering F1 score falling from
45.6 to 20.6 when transferring from a 4-billion-parameter model to a
14-billion-parameter model, while some transfers in the opposite direction
stay close to the target's own results \cite{Yun2026HeteroFold}. A cache
with the right shape and a passing compatibility check is therefore not
necessarily good enough, and a predictor that scores well on the
conversations it was fitted on may not do so on new ones.

\paragraph{H9, altered cache.} Retained KV state can change a model's
behaviour even when the text that produced it is no longer part of the
request \cite{Zhang2026Aborted}. A cache that has been corrupted or altered
could therefore steer the response without any visible change to the
prompt.

\paragraph{H10, timing side channel.} Shared prefix caches in LLM serving
systems have been shown to leak information through timing: a cache hit
returns its first token sooner, which lets an observer infer what other
users have sent \cite{Song2024EarlyBird}. \progname{} adds another fast
path, so the same risk applies.

\section{Failure Mode and Effect Analysis}
\label{sec:fmea}

Table~\ref{tab:fmea} analyzes how each component can fail. Each row gives
the recommended action and the \srs{} requirement that implements it.

\begin{landscape}
\footnotesize
\setlength{\tabcolsep}{3pt}
\begin{longtable}{@{}p{0.04\linewidth}p{0.10\linewidth}p{0.13\linewidth}%
p{0.13\linewidth}p{0.13\linewidth}p{0.11\linewidth}p{0.16\linewidth}%
p{0.07\linewidth}p{0.05\linewidth}@{}}
\caption{Failure Mode and Effect Analysis}\label{tab:fmea}\\
\toprule
ID & Component & Failure mode & Effects & Causes & Detection &
Recommended action & Req. & Hazard\\
\midrule\endfirsthead
\toprule
ID & Component & Failure mode & Effects & Causes & Detection &
Recommended action & Req. & Hazard\\
\midrule\endhead
F1 & Request validation & Accepts a malformed or oversized request &
Out-of-memory error or wrong output; other requests slowed & Missing size
or setting checks & Rejection tests with invalid inputs (T-ADM) & Reject
before allocating resources; check sizes against the configuration & FR-01,
SR-07 & H6\\
F2 & Compatibility checker & Accepts a cache predicted for a different
model, tokenizer, format or position layout & Garbled or subtly wrong
answer & Missing metadata treated as a pass; equal cache sizes treated as a
match & Each check forced to fail separately (T-COMP, T-MATCH) & Treat
missing information as a failed check; require every check to pass & FR-03,
FR-04, SR-02 & H1\\
F3 & Compatibility checker & Accepts a cache from before a model or
configuration change & Wrong answer & Version not recorded; restart without
invalidation & Change the model or configuration, then reuse an old cache
(T-CONFIG) & Record versions with each cache; invalidate on restart or
change & SR-05, FR-13 & H3, H13\\
F4 & KV cache predictor & A tested, compatible predicted cache lowers
answer quality & Plausible but wrong answers, worst when predicting a
larger model from a smaller one & Approximate prediction; conversation
unlike the fitting data & Held-out prediction error and output tests
(T-PRED, T-QUAL); run-time check on the last history tokens & Approve
each setup only within a quality-loss limit; at run time, recompute the
last few history tokens on top of the predicted cache and fall back when
they disagree & SR-09, SR-10, FR-14, FR-16 & H8\\
F5 & KV cache predictor & Predicted cache is corrupted before loading &
Steered or wrong answer & Storage error; software fault; untrusted source
in a future release & Integrity check fails & Record a checksum and the
predictor version with each cache; reject on mismatch & SR-11, SR-02 &
H9\\
F6 & KV cache predictor & Tokens, layers or heads aligned wrongly between
source and target & Keys and values placed at the wrong positions; garbled
answer & Tokenizers split text differently; alignment rules not updated
for a new model pair & Alignment checks in T-PRED; run-time check on the
last history tokens & Record alignment with each cache; reject a cache whose
alignment does not match the configured rules & FR-16, SR-02, SR-10 & H1,
H8\\
F7 & KV cache predictor & Produces invalid tensors (wrong shape,
non-finite values) & Request crashes or produces nonsense & Numerical
overflow; input outside the fitted range & Forced invalid outputs
(T-FAIL) & Validate shape and values before loading; fall back to native
prefill & FR-07, FR-16 & H1, H6\\
F8 & KV cache predictor & Prediction and loading take longer than native
prefill & Slower first token than without \progname{} & Short shared
history; large predictor; slow loading path & TTFT benchmarks against
native prefill (T-PERF) & Use prediction only where tests show a gain;
otherwise use native prefill & NFR-04, NFR-08 & H15\\
F9 & Compatibility checker & Reuses a cache from another request's sharing
group & Private history exposed & Sharing group not checked; request names
its own group & Cross-group isolation tests (T-ISOLATION) & Sharing group
assigned by the host; permission check required & SR-01, SR-04 & H2\\
F10 & KV loader & Reuse is measurably faster than normal processing & An
observer infers other users' conversations & Caches reused across users &
Review of sharing settings; timing tests & Reuse only within one sharing
group; disable cross-user sharing by default; document the remaining risk &
SR-12 & H10\\
F11 & KV loader & Cache evicted or overwritten while in use & Corrupted
output for an active request & Reference-counting fault; eviction policy &
Repeated lifecycle tests (T-LIFE) & Pin caches while any request uses them
& SR-03, FR-09 & H4\\
F12 & KV loader & Loading fails after output has started & Partial answer
& GPU error; memory shortage & Forced failure tests (T-FAIL) & End with a
failure status; never append a restarted answer & FR-07, SR-08 & H7\\
F13 & Native prefill fallback & Fallback is not triggered or fails & A
valid request fails & Logic fault; fallback disabled & Missing-cache and
unsupported-pair tests (T-NATIVE) & Fall back on any failed check; keep the
normal path always available & FR-06, NFR-06 & H1, H6\\
F14 & Native prefill fallback & Many requests fall back at the same time &
Delay spike and timeouts for all requests & A configuration change or
predictor fault invalidates many caches at once & Fallback-rate and queue
metrics & Admit prefills only within the engine's budget; queue a bounded
number and reject the rest with a retry status; turn reuse off when the
fallback rate exceeds a limit & SR-13, NFR-08 & H14\\
F15 & KV cache predictor & Predicted caches are never used & Memory
pressure, evictions and wasted GPU time & Cache predicted for a model the
router did not select & Memory and prediction metrics & Cap the memory for
predicted caches and expire unused ones & SR-14 & H12, H6\\
F16 & Request lifecycle & Cancelled or timed-out request keeps its memory
& Memory leak leading to out-of-memory errors & Missing release path &
Repeated cancellation tests (T-LIFE, T-CANCEL) & Release resources on every
final status & FR-09, FR-15 & H6\\
F17 & Output streaming & Tokens attached to the wrong request & A user
receives another user's answer & Request IDs mixed under concurrency &
Overlapping-request tests (T-STATE) & Keep one ordered output stream and one
final status per request & FR-08 & H2\\
F18 & Logging and metrics & Raw prompts, answers or cache data written to
logs & Private data exposed & Debug logging left on & Scan of logs after
tests (T-PRIVACY) & Log only IDs and status or reason codes & SR-06 & H5\\
F19 & Benchmark tools & Baseline differs from the \progname{} run, or
failures and preparation cost are left out & Overstated benefit & Different
settings; selective reporting & Configuration comparison (T-EVAL, T-REPORT)
& Run matched workloads; report failures, runs with no gain, and cache
preparation cost & FR-11, FR-12 & H11\\
F20 & Configuration registry & An untested setup is enabled & Wrong or
low-quality answers & Operator error & Attempt to enable without evidence
(T-QUAL) & Keep reuse off for setups without passing test records & FR-13,
FR-14 & H1, H8\\
F21 & Benchmark tools & Predictor evaluated on conversations it was
fitted on & Reported accuracy higher than in real use & Data split done
after fitting, or overlapping conversations & Check of the recorded data
split (T-QUAL, T-REPORT) & Split conversations into fitting and held-out
groups before fitting; report the split & FR-14, FR-12 & H11, H8\\
\bottomrule
\end{longtable}
\end{landscape}

\section{Safety and Security Requirements}
\label{sec:requirements}

The existing \srs{} requirements SR-01 to SR-08 mitigate H1--H7 and are
referenced in Table~\ref{tab:fmea}. The analysis found the following new
requirements, which should be added to the \srs{} security and robustness
sections. Their fit criteria use the constants in
Table~\ref{tab:constants}.

\begin{description}[style=nextline]
  \item[SR-09 --- Limit the quality loss of each cross-model setup (H8)]
  \progname{} shall enable a cross-model setup only if, on its declared test
  set, which shares no conversations with the data used to fit the
  predictor, the selected model's task scores with the predicted cache are
  no more than \texttt{QUALITY\_LOSS\_LIMIT} below its scores with native
  prefill. \emph{Fit criterion:} T-QUAL compares both runs on the same test
  set; a setup above the limit stays disabled.
  \item[SR-10 --- Check predicted caches at run time (H8)] For each reuse,
  \progname{} shall recompute the last \texttt{CANARY\_TOKENS} tokens of
  the shared history with the target model, using the predicted cache for
  all earlier positions, and compare the recomputed keys and values with
  the predicted ones at those positions. If the difference exceeds
  \texttt{CANARY\_LIMIT}, it shall use native prefill. The check costs
  about \texttt{CANARY\_TOKENS} tokens of prefill, not the whole history.
  It detects badly wrong predictions but does not prove a cache correct;
  SR-09 provides the comparison with native prefill. \emph{Fit criterion:}
  a test injects a degraded cache and checks that the request falls back
  and records the reason; T-PERF shows that the check keeps time to first
  token below native prefill for supported setups.
  \item[SR-11 --- Check cache integrity (H9)] Each predicted cache shall
  carry a checksum and a record of its source model and predictor version;
  \progname{} shall reject a cache whose checksum or record does not
  match. \emph{Fit
  criterion:} a test alters one byte of a cache and checks that it is
  rejected.
  \item[SR-12 --- Limit timing leaks (H10)] \progname{} shall reuse a cache
  only within the sharing group it was predicted for, with reuse across users
  disabled by default; the documentation shall state the remaining timing
  risk. \emph{Fit criterion:} T-ISOLATION shows no reuse across groups; a
  documentation review finds the risk statement.
  \item[SR-13 --- Limit fallback load (H14)] \progname{} shall admit
  fallback prefills only within the engine's configured prefill budget,
  queue at most \texttt{QUEUE\_LIMIT} further requests, and reject any
  beyond that with a status telling the caller to retry, rather than
  letting them time out. It shall also turn cross-model reuse off
  automatically when the fallback rate over the last
  \texttt{FALLBACK\_WINDOW} requests exceeds \texttt{FALLBACK\_RATE\_LIMIT},
  so no prediction work is wasted, and report this to the operator.
  \emph{Fit criterion:} a test invalidates all caches at once under load;
  admitted requests stay within \texttt{REGRESSION\_LIMIT} of the same load
  without \progname{} (NFR-08), excess requests receive the retry status
  instead of timing out, and reuse is turned off.
  \item[SR-14 --- Bound memory for predicted caches (H12)] Memory held by
  predicted but unused caches shall not exceed
  \texttt{PREDICTED\_CACHE\_LIMIT}, and unused caches shall expire after
  \texttt{PREDICTED\_CACHE\_TTL}. \emph{Fit criterion:} a test predicts
  caches that are never used and checks that memory stays within the limit.
\end{description}

\begin{table}[h]
\centering
\caption{New constants (values to be set by the team)}
\label{tab:constants}
\begin{tabularx}{\textwidth}{@{}lX@{}}
\toprule Constant & Meaning\\\midrule
\texttt{QUALITY\_LOSS\_LIMIT} & Largest accepted drop in task score for a
cross-model setup. \todo{choose a value with the team and advisor}\\
\texttt{CANARY\_TOKENS} & Number of final history tokens recomputed for
the run-time check. \todo{choose a value}\\
\texttt{CANARY\_LIMIT} & Largest accepted difference between the
recomputed and predicted keys and values for those tokens.
\todo{choose a value}\\
\texttt{FALLBACK\_WINDOW} & Number of recent requests used to compute the
fallback rate. \todo{choose a value}\\
\texttt{FALLBACK\_RATE\_LIMIT} & Fallback rate above which reuse is turned
off. \todo{choose a value}\\
\texttt{QUEUE\_LIMIT} & Number of requests that may wait for prefill
before further requests are rejected. \todo{choose a value}\\
\texttt{PREDICTED\_CACHE\_LIMIT} & Memory allowed for predicted but unused
caches. \todo{choose a value}\\
\texttt{PREDICTED\_CACHE\_TTL} & Time after which an unused predicted cache
expires. \todo{choose a value}\\
\bottomrule
\end{tabularx}
\end{table}

\paragraph{Changes needed in the \srs{}.}
\begin{itemize}
  \item Add SR-09 to SR-14, their constants, and H8--H15 to the traceability
  table, which resolves open issue OI6. FR-16 should also trace to H8, H13
  and H15, not only H1 and H3.
  \item \textbf{NFR-05 conflicts with H8.} NFR-05 requires every reuse
  setup to match the native scores within \texttt{SCORE\_ATOL} and to
  produce identical greedy tokens. The \srs{} itself states that cache
  prediction may be approximate, so no cross-model setup could pass, and
  FR-14 would never enable reuse. NFR-05
  should keep the exact check for same-model reuse and use SR-09 and SR-10
  for cross-model reuse. \todo{agree on this change with the SRS author}
  \item Note in A1 that H9 is still analyzed (Section~\ref{sec:assumptions}).
\end{itemize}

\section{Roadmap}
\label{sec:roadmap}

Table~\ref{tab:roadmap} shows which requirements are planned within the
capstone and which are left for later. Requirements for the capstone are
the ones needed for a correct, safe single-server demonstration.

\begin{table}[h]
\centering
\caption{Roadmap for the safety and security requirements}
\label{tab:roadmap}
\begin{tabularx}{\textwidth}{@{}p{0.27\textwidth}p{0.16\textwidth}X@{}}
\toprule Requirements & When & Reason\\\midrule
SR-01 to SR-08 & Capstone & Needed for any correct reuse; already in the
\srs{}.\\
SR-09, SR-10 & Capstone & Cross-model prediction is approximate (H8), so
quality checks are part of the core product.\\
SR-11 (checksum and source) & Capstone & Cheap to build and detects
corruption.\\
SR-12 (sharing-group limit) & Capstone & Configuration and documentation
only.\\
SR-13, SR-14 & Capstone, after the core path works & Needed for load tests,
but depend on the reuse path.\\
Signed caches and caches uploaded by users & Future & A1 excludes untrusted
sources in the first release; accepting them would need authentication of
the cache source.\\
Timing side-channel defences, such as padding response times & Future &
Costs part of the speedup; SR-12 reduces the risk for now.\\
Hazards of operation across several servers & Future & Outside the
\srs{} scope.\\
\bottomrule
\end{tabularx}
\end{table}

\bibliographystyle{plain}
\bibliography{../../refs/References}

\newpage{}

\section{Appendix --- Generative AI Use Disclosure}

\subsection{AI Tools Used}

Claude Code, running Claude Opus 5.5 (Anthropic).

\subsection{How and Where Used}

\begin{itemize}
  \item \textbf{Drafting:} Claude drafted all sections of this document from
  the team's SRS draft, the course Hazard Analysis lecture, the template and
  the Hazard Analysis checklist. This includes the hazard list, the FMEA
  table, SR-09 to SR-14 and the roadmap.
  \item \textbf{Consistency checking:} Claude compared the SRS requirement
  IDs and draft hazard labels with this document, and identified the
  conflict between NFR-05 and approximate cross-model prediction.
  \item \textbf{Literature:} Claude found the three papers cited for H8,
  H9 and H10, and checked their titles, authors and the reported numbers on
  the arXiv pages and the HeteroFold repository.
\end{itemize}

\subsection{Verification of AI-Generated Content}

The citations were checked against the arXiv abstract pages, and the
H8 figures against the results table published with the HeteroFold code.
The requirement IDs were checked against the SRS draft on the \texttt{SRS}
branch. The team reviewed every hazard, FMEA row and new requirement
before submission. The team is
responsible for the accuracy, completeness and consistency of this
document.

\newpage{}

\section{Appendix --- Reflection}

\input{../Reflection.text}

\begin{enumerate}
    \item What went well while writing this deliverable?

    After meeting with our supervisor, the team had a clear idea of the
    project's direction. Identifying the hazards and losses went well
    because, by this point, we also understood the technology well from
    reviewing the state-of-the-art papers on KV cache reuse and cross-model
    transfer, so we could see where the system could fail and who would be
    affected.

    \item What pain points did you experience during this deliverable, and how
    did you resolve them?

    The hardest part was coordinating between the writer of the SRS and the
    writer of this Hazard Analysis, who worked at the same time and had to
    agree on the functional, non-functional and security requirements. We
    resolved this by using the SRS identifiers throughout this document and
    collecting every change the SRS needs, including the new requirements
    and the NFR-05 conflict, in Section~\ref{sec:requirements} so that the SRS
    author could review and apply them in one place.

    \item Which of your listed risks had your team thought of before this
    deliverable, and which did you think of while doing this deliverable? For
    the latter ones (ones you thought of while doing the Hazard Analysis), how
    did they come about?

    H1--H7 were identified before this
    deliverable, while writing the SRS. They came from discussions within
    the team and with our supervisor, and from research papers on KV cache
    reuse, which showed when a cache can safely be reused. H8--H15 came up
    during this analysis. H8 came from reading published cross-model
    transfer results, which showed large and direction-dependent quality
    loss. H10 came from work on timing leaks in prefix caching. H11 to H14
    came from asking what happens when many requests fail at once, or when
    predicted caches are never used. H15 came from the SRS assumption that a
    predictor cheaper than prefill is still an open question.

    \item Other than the risk of physical harm (some projects may not have any
    appreciable risks of this form), list at least 2 other types of risk in
    software products. Why are they important to consider?

    Privacy risk: software can reveal
    information users expected to stay private, as in H2, H5 and H10, and
    such a leak cannot be undone. Risk of misleading results: software that
    produces confident but wrong output, as in H8 and H11, can lead people
    to make poor decisions, because nothing tells them the output is
    unreliable.
\end{enumerate}

\end{document}
